\subsection{Merge sort}
\label{sec:merge_sort}

Merge sort resuelve el \hyperref[sec:problema]{problema de ordenamiento}. La idea es:
\begin{enumerate}
    \item Dividir el arreglo en dos mitades.
    \item Conquistar: Ordenar cada mitad recursivamente (i.e., dividir cada mitad otra vez en la mitad, hasta que queden segmentos de un elemento, que entonces ya están ordenados)
    \item Combinar los subarreglos ordenados.
\end{enumerate}

\subsubsection{Dividir y conquistar}

Dividir no tiene ciencia: se calcula el índice de la mitad del arreglo para determinar el índices de inicio y fin de cada segmento.

Conquistar tampoco: se hace el llamado recursivo a cada mitad. Al partir en mitades eventualmente se obtienen subarreglos de longitud uno, ya ordenados.

Por ahora no se piensa en cómo combinar, se encapsula ese paso en una función \inline{merge} que se implementará después. Con eso, este es el seudocódigo:
\begin{pseudocode}
def merge_sort(nums: arreglo[int], lo: int, hi: int):
    if lo < hi:
        |\bxl{mid}|mid = (hi - lo) // 2 + lo|\bxr{mid}|

        merge_sort(nums, lo, mid)
        merge_sort(nums, mid + 1, hi)

        merge(nums, lo, mid, hi)
\end{pseudocode}
\begin{annotations}
    \codebox[xshift=0.5cm,width=7cm]{mid}{\annMitadSinOverflow{lo}{hi}}
\end{annotations}

Los índices \inline{lo} y \inline{hi} como parámetros son necesarios porque indican qué segmento del arreglo se está procesando. La invocación inicial, que resuelve el problema completo, usa los índices extremos del arreglo: \inline{merge_sort(nums, 0, len(nums)-1)}. Nótese que el arreglo realmente no se divide en cada llamado recursivo, sino que se delimitan distintas regiones mediante índices. Eso permite que el algoritmo sea \concept{\textit{in-place}}, es decir, modifique directamente el arreglo original en lugar de retornar uno nuevo.

En cada llamado recursivo ocurre lo siguiente:
\begin{enumerate}
    \item Si el segmento tiene longitud 1 (o 0), se alcanza el caso base: \inline{lo} no es menor que \inline{hi} y la función termina.
    
    \item Si no, el segmento se divide conceptualmente en dos mitades y se realizan llamados recursivos para cada una. El llamado a la mitad derecha empieza \emph{después} de que termina el llamado a la mitad izquierda.
    
    \item Una vez ambos llamados retornan, se ejecuta \inline{merge}, que mezcla ambas mitades ordenadas in-place dentro del rango delimitado por \inline{lo} y \inline{hi}.
\end{enumerate}

Para arreglos grandes, el caso base solo se alcanza después de muchos llamados recursivos, cuando el arreglo se ha dividido lo suficiente. Ergo, en las primeras invocaciones, siempre ocurre que el segmento se divide conceptualmente: se apilan llamados a \inline{merge_sort}, sin retornar ni combinar.

Cuando algún llamado alcanza el caso base, vuelva a la función que la invocó justo después de donde la invocó (al final de la línea 5 o de la línea 6) y pasa una de dos cosas:
\begin{itemize}
    \item Se realiza otro llamado recursivo (que pronto alcanzará el caso base, porque ya el segmento es muy pequeño)
    \item Se ejecuta por primera vez el paso de combinar, mezclando esos dos segmentos (que son pequeños, porque hace poco se llegó al caso base)
\end{itemize}
Después de combinar, la función termina y retorna a quien la había invocado. Así, los llamados a \inline{merge_sort} que se habían acumulado empiezan a terminar, uno tras otro: la recursión comienza a \concept{desenrollarse}.

\subsubsection{Combinar las mitades}

Ahora sí: cómo combinar las mitades.

La idea es intuitiva: si cada mitad fuera una pila de cartas, el método natural sería comparar las cartas en el tope de cada pila y tomar la menor cada vez. Eso solo funciona porque cada mitad está ordenada, entonces las cartas que se van tomando van quedando ordenadas.

El seudocódigo es el siguiente:
\begin{pseudocode}
def merge(nums: arreglo[int], lo: int, mid: int, hi: int):
    |\br{cpi}|left = |\textrm{copia de}| nums |\textrm{desde}| lo |\textrm{hasta}| mid
    right = |\textrm{copia de}| nums |\textrm{desde}| mid+1 |\textrm{hasta}| hi|\br{cpf}|

    |\bxl{ir}|l = 0|\bxr{ir}|
    r = 0

    |\br{ai}|while l < len(left) and r < len(right):
        if left[l] |\bxl{es}|<=|\bxr{es}| right[r]:
            nums[lo + l + r] = left[l]
            l += 1
        else:
            nums[lo + l + r] = right[r]
            r += 1|\br{af}|

    |\br{ii}|while l < len(left):
        nums[lo + l + r] = left[l]
        l += 1|\br{if}|

    |\br{di}|while r < len(right):
        nums[lo + l + r] = right[r]
        r += 1|\br{df}|
\end{pseudocode}
\begin{annotations}
    \codebrace[width=5.5cm]{cpi}{cpf}{Se crean arreglos nuevos con cada mitad. La complejidad espacial es entonces \(\Theta(n)\).}
    \codebox[xshift=0.5cm]{ir}{\inline{l} recorre  \inline{left}}
    \codebox[notepos={(cbox@es.east)+(2cm,0.8cm)}, curve={bend left=25}, width=10cm]{es}{\texttt{<=} en lugar de \texttt{<} al comparar los elementos de cada mitad para que haya estabilidad.}
    \codebrace[xshift=4.5cm, width=6cm]{ai}{af}{Lógica principal, cuando ambos subarreglos tienen elementos: se toma el menor y se añade a \inline{nums}. \vspace{0.1cm}\newline \inline{lo+l+r} es el índice de \inline{nums} en el que se está: \inline{lo} donde inicia el segmento, \inline{l} los elementos de \inline{left} que se han agregado y \inline{r} los de \inline{right}}
    \codebrace[xshift=3.7cm, width=7.5cm]{ii}{if}{Si quedan elementos en el arreglo de la izquierda, ya no quedan en el de la derecha. Se añade el excedente de la izquierda a \inline{nums}.}
    \codebrace[xshift=3.7cm, width=6.5cm]{di}{df}{Se añade el excedente de la derecha a \inline{nums}. Si este while corre, es porque el anterior no corrió y viceversa.}
\end{annotations}

Ese algoritmo es \concept{estable}: mantiene el orden relativo de los elementos iguales. O sea, en caso de empate, deja primero al que estaba primero en el arreglo original. En la línea del comentario sobre estabilidad, se elige el elemento de la izquierda en caso de empate para que continúe ubicado antes que el de la derecha en la solución final.

\subsubsection{Complejidad}
\label{sec:merge_sort_complejidad}

Una forma de razonar sobre la complejidad es analizando el \hyperref[sec:arbol_de_ejecucion]{árbol de ejecución}:
\begin{center}
\begin{forest}
  for tree={
    draw, rounded corners,
    font=\ttfamily\small,
    inner sep=4pt,
    l sep=1.5em,
    s sep=0.6em,
    edge={-latex},
  }
  [{8}
    [{4}
      [{2}
        [{1}]
        [{1}]
      ]
      [{2}
        [{1}]
        [{1}]
      ]
    ]
    [{4}
      [{2}
        [{1}]
        [{1}]
      ]
      [{2}
        [{1}]
        [{1}]
      ]
    ]
  ]
\end{forest}
\end{center}

Cada nodo está etiquetado con el tamaño del segmento que procesa. Como cada llamado divide su segmento exactamente por la mitad, el árbol es balanceado y se puede razonar nivel por nivel en lugar de nodo por nodo:
\begin{itemize}
  \item En el nivel 0 (la raíz) hay un solo llamado, sobre un segmento de tamaño \(n\); su costo no recursivo, el costo de combinar, es \(\Theta(n)\).
  \item En el nivel 1 hay dos llamados, cada uno sobre un segmento de tamaño \(n/2\). El costo de combinar cada uno es de \(\Theta(n/2)\), para un total de \(2 \cdot \Theta(n/2) = \Theta(n)\).
  \item En general, en el nivel \(i\) hay \(2^i\) llamados, cada uno sobre un segmento de tamaño \(n/2^i\) cuyo costo de combinación es es \(\Theta(n/2^i)\), para un costo total de \(\Theta(n)\).
\end{itemize}
Es decir, la suma de los costos no recursivos para cada nivel es de \(\Theta(n)\), independientemente del nivel.

Con eso, solo hace falta contar los niveles, la altura \(h\) del árbol, para saber cuántas veces se paga ese costo. Como la raíz representa un segmento de tamaño \(n\) que en cada nivel se divide a la mitad, solo hace falta saber cuántas veces se debe dividir \(n\) entre \(2\) hasta llegar a \(1\).
\begin{gather*}
  n \underbrace{{} \cdot \frac{1}{2} \cdot \frac{1}{2} \cdots \frac{1}{2}}_{h \ \text{veces}} = 1 \\
  \frac{n}{2^h} = 1 \\
  2^h = n \\
  h = \log_2(n)
\end{gather*}
Es decir, se paga un costo lineal por cada nivel y la cantidad de niveles es logarítmica en función de la entrada, por lo cual la complejidad temporal de merge sort es \(\Theta(n \log n)\).

Otra forma de calcular la complejidad es mediante la \hyperref[sec:complejidad-temporal-recursion]{ecuación de recurrencia}: en el caso base el algoritmo simplemente retorna, con costo constante; en el otro caso, el costo no recursivo es el de combinar, \(\Theta(n)\), sumado a los dos llamados recursivos, cada uno sobre una mitad del arreglo de tamaño aproximadamente \(n/2\):
\[
T(n) =
\begin{cases}
  \Theta(1) & n \le 1. \\
  \Theta(n) + 2\,T(n/2) & n > 1.
\end{cases}
\]
La solución de \(T(n) = \Theta(n) + 2\,T(n/2)\) es precisamente \(\Theta(n \log n)\) (véase el apéndice \ref{apx:ecuaciones_recursivas}).

\subsubsection{Optimalidad}

Merge sort es un \concept{algoritmo de ordenamiento basado en comparaciones}: la única forma en la que decide el orden relativo de dos elementos es comparándolos entre sí. Se dice que merge sort es \concept{asintóticamente óptimo} porque ningún algoritmo de ese tipo puede tener, en el peor caso, un costo menor a \(\Omega(n \log n)\).

% Revisar de aquí para abajo
Todo algoritmo de ordenamiento basado en comparaciones se puede representar mediante un \concept{árbol de decisión}: un árbol binario en el que cada nodo interno representa una comparación entre dos elementos, sus dos hijos representan los dos posibles resultados de esa comparación, y cada hoja representa una permutación de la entrada que el algoritmo podría retornar como resultado.

El algoritmo debe funcionar correctamente sin importar cuál de las \(n!\) permutaciones posibles de tamaño \(n\) sea la entrada. Por ende, el árbol de decisión necesita al menos una hoja por cada una de esas \(n!\) permutaciones: si le faltara alguna, existiría una entrada para la cual el algoritmo jamás llegaría a la salida correcta.

La altura del árbol es el número de comparaciones que hace el algoritmo en el peor caso, pues cada camino de la raíz a una hoja corresponde a una posible ejecución. Como el árbol es binario, uno de altura \(h\) tiene a lo sumo \(2^h\) hojas, así que:
\[2^h \ge n! \implies h \ge \log_2(n!)\]
Acotando ese logaritmo:
\[
\log_2(n!) = \sum_{k=1}^n \log_2 k \ge \sum_{k=n/2}^{n} \log_2 k \ge \frac{n}{2} \log_2 \frac{n}{2} \in \Omega(n \log n)
\]
Ergo, todo algoritmo de ordenamiento basado en comparaciones necesita, en el peor caso, \(\Omega(n \log n)\) comparaciones. Como \inline{merge_sort} logra \(O(n \log n)\), es asintóticamente óptimo.

\practicar[leetcode]{Sort an Array}{https://leetcode.com/problems/sort-an-array/}